\documentclass[11pt]{article}
\usepackage{../../latex/stat4ml_style}

\title{Chapter 6: Markov Decision Processes and Dynamic Programming \\ \large Part III --- Statistical Foundations of Deep Learning \& Reinforcement Learning}
\author{Omid Safarzadeh}
\date{}

\begin{document}
\maketitle

\begin{abstract}
\noindent
Reinforcement learning begins not with an algorithm but with a fixed point. A
Markov decision process turns ``act well over time'' into a precise object --- the
value function, the expected discounted reward from each state --- and the value
function is characterised by the Bellman equation, which says a state's value is
its immediate reward plus the discounted value of where it leads. The reason this
equation can be \emph{solved} is one clean fact: the Bellman operator is a
contraction, shrinking distances by the discount factor $\gamma$, so by Banach's
theorem it has a unique fixed point that repeated application converges to. This
chapter proves the contraction, reads value and policy iteration off it, and ---
when the transition model is unknown --- estimates the value from sampled
experience by Monte Carlo and temporal-difference learning, the latter being the
stochastic approximation of Chapter~3 in disguise.
\end{abstract}

\tableofcontents

% ======================================================================
\section{Markov decision processes and the value function}
% ======================================================================

\begin{definition}[MDP]
A \keyterm{Markov decision process} is $(\mathcal{S},\mathcal{A},P,r,\gamma)$:
states $\mathcal{S}$, actions $\mathcal{A}$, a transition kernel
$P(s'\mid s,a)$, a reward $r(s,a)$, and a \keyterm{discount} $\gamma\in[0,1)$. A
\keyterm{policy} $\pi(a\mid s)$ chooses actions; its \keyterm{value function} is
the expected discounted return
\[
  V^\pi(s) = \mathbb{E}_\pi\!\left[\sum_{t=0}^\infty \gamma^t r(s_t,a_t)\,\Big|\, s_0=s\right].
\]
\end{definition}

The sum telescopes into a one-step relation, because the return from $s$ is the
reward now plus $\gamma$ times the return from the next state.

\begin{definition}[Bellman equations]
The value of a fixed policy satisfies the \keyterm{Bellman expectation equation}
\[
  V^\pi(s) = \sum_a \pi(a\mid s)\Big[r(s,a) + \gamma\sum_{s'}P(s'\mid s,a)V^\pi(s')\Big],
\]
and the optimal value the \keyterm{Bellman optimality equation}
\[
  V^*(s) = \max_a\Big[r(s,a) + \gamma\sum_{s'}P(s'\mid s,a)V^*(s')\Big].
\]
\end{definition}

The left panel of the figure shows an optimal value function $V^*$ on a
slippery gridworld, together with the greedy policy it induces --- each arrow
points to the action that maximises the bracket above.

% ======================================================================
\section{The Bellman operator is a contraction}
% ======================================================================

Write the optimality equation as a map on value functions: the \keyterm{Bellman
optimality operator} $T$ sends a value function $V$ to
\[
  (TV)(s) = \max_a\Big[r(s,a) + \gamma\sum_{s'}P(s'\mid s,a)V(s')\Big].
\]
$V^*$ is exactly a fixed point of $T$. The following is why that fixed point
exists, is unique, and can be found by iterating.

\begin{theorem}[$T$ is a $\gamma$-contraction in the sup norm]
For any value functions $V,U$, $\;\|TV - TU\|_\infty \le \gamma\,\|V-U\|_\infty$.
\end{theorem}
\begin{proof}
Fix a state $s$. Using $|\max_a f(a)-\max_a g(a)|\le \max_a|f(a)-g(a)|$,
\[
  |(TV)(s)-(TU)(s)|
  \le \max_a \gamma\sum_{s'}P(s'\mid s,a)\,|V(s')-U(s')|
  \le \gamma\,\|V-U\|_\infty \max_a\sum_{s'}P(s'\mid s,a).
\]
The rewards cancel, and $\sum_{s'}P(s'\mid s,a)=1$, so the right side is
$\gamma\|V-U\|_\infty$. Taking the max over $s$ gives the claim.
\end{proof}

\begin{theorem}[Value iteration converges]
$T$ has a unique fixed point $V^*$, and $V_{k+1}=TV_k$ satisfies
$\|V_k-V^*\|_\infty \le \gamma^k\|V_0-V^*\|_\infty \to 0$.
\end{theorem}
\begin{proof}
A contraction on the complete metric space $(\mathbb{R}^{|\mathcal{S}|},
\|\cdot\|_\infty)$ has a unique fixed point and its iterates converge to it
(Banach). Applying the contraction $k$ times to $V_0$ and $V^*=TV^*$ gives
$\|V_k-V^*\|_\infty=\|T^kV_0 - T^kV^*\|_\infty\le\gamma^k\|V_0-V^*\|_\infty$.
\end{proof}

The middle panel is this bound in action: the sup-norm error falls geometrically,
staying under the $\gamma^k$ line. \keyterm{Policy iteration} is the other route
to $V^*$ --- alternate evaluating the current policy (solving the linear Bellman
expectation equation) and acting greedily with respect to that value --- and it
converges in finitely many steps, because there are finitely many policies and
each step strictly improves.

% ======================================================================
\section{When the model is unknown: Monte Carlo and TD}
% ======================================================================

Dynamic programming needs $P$ and $r$. Reinforcement learning usually does not
have them, only sampled experience, so it \emph{estimates} the value. Two ways,
and they are the two integral-estimation styles of Chapter~4 applied to the
return. \keyterm{Monte Carlo} evaluation waits until an episode ends and moves
each visited state's estimate toward the observed return $G$:
\[
  V(s) \leftarrow V(s) + \alpha\big(G - V(s)\big).
\]
\keyterm{Temporal-difference} learning does not wait; it bootstraps, moving toward
the one-step estimate $r + \gamma V(s')$ immediately:
\[
  V(s) \leftarrow V(s) + \alpha\big(\,\underbrace{r + \gamma V(s') - V(s)}_{\text{TD error}}\,\big).
\]

TD(0) is Robbins--Monro stochastic approximation (Chapter~3) on the Bellman
expectation equation: the TD error is a noisy, unbiased sample of ``how wrong the
current value is,'' and the step-size conditions $\sum\alpha_t=\infty$,
$\sum\alpha_t^2<\infty$ are exactly what guarantee convergence. Because TD updates
every step and uses its own value estimate as a target, it typically has lower
variance than Monte Carlo --- the right panel, where TD(0) reaches a given
accuracy in far fewer episodes than constant-$\alpha$ Monte Carlo on the same
random walk.

\begin{summarybox}
An MDP $(\mathcal{S},\mathcal{A},P,r,\gamma)$ defines the value function
$V^\pi$, characterised by the Bellman equations; the optimal value $V^*$ is the
fixed point of the Bellman optimality operator $T$. $T$ is a $\gamma$-contraction
in the sup norm, so by Banach it has a unique fixed point and value iteration
converges geometrically, $\|V_k-V^*\|_\infty\le\gamma^k\|V_0-V^*\|_\infty$;
policy iteration converges in finitely many strict-improvement steps. Without a
model, the value is estimated from experience: Monte Carlo moves toward the
realised return, TD(0) bootstraps toward $r+\gamma V(s')$ --- the latter is the
Chapter~3 stochastic approximation of the Bellman equation, lower-variance than
Monte Carlo, and the foundation for the policy-gradient methods of Chapter~7.
\end{summarybox}

\end{document}
